% Options for packages loaded elsewhere
\PassOptionsToPackage{unicode}{hyperref}
\PassOptionsToPackage{hyphens}{url}
\PassOptionsToPackage{dvipsnames,svgnames,x11names}{xcolor}
\documentclass[
  11pt,
]{article}
\usepackage{xcolor}
\usepackage[margin=1in]{geometry}
\usepackage{amsmath,amssymb}
\setcounter{secnumdepth}{-\maxdimen} % remove section numbering
\usepackage{iftex}
\ifPDFTeX
  \usepackage[T1]{fontenc}
  \usepackage[utf8]{inputenc}
  \usepackage{textcomp} % provide euro and other symbols
\else % if luatex or xetex
  \usepackage{unicode-math} % this also loads fontspec
  \defaultfontfeatures{Scale=MatchLowercase}
  \defaultfontfeatures[\rmfamily]{Ligatures=TeX,Scale=1}
\fi
\usepackage{lmodern}
\ifPDFTeX\else
  % xetex/luatex font selection
\fi
% Use upquote if available, for straight quotes in verbatim environments
\IfFileExists{upquote.sty}{\usepackage{upquote}}{}
\IfFileExists{microtype.sty}{% use microtype if available
  \usepackage[]{microtype}
  \UseMicrotypeSet[protrusion]{basicmath} % disable protrusion for tt fonts
}{}
\makeatletter
\@ifundefined{KOMAClassName}{% if non-KOMA class
  \IfFileExists{parskip.sty}{%
    \usepackage{parskip}
  }{% else
    \setlength{\parindent}{0pt}
    \setlength{\parskip}{6pt plus 2pt minus 1pt}}
}{% if KOMA class
  \KOMAoptions{parskip=half}}
\makeatother
\usepackage{color}
\usepackage{fancyvrb}
\newcommand{\VerbBar}{|}
\newcommand{\VERB}{\Verb[commandchars=\\\{\}]}
\DefineVerbatimEnvironment{Highlighting}{Verbatim}{commandchars=\\\{\}}
% Add ',fontsize=\small' for more characters per line
\newenvironment{Shaded}{}{}
\newcommand{\AlertTok}[1]{\textcolor[rgb]{1.00,0.00,0.00}{\textbf{#1}}}
\newcommand{\AnnotationTok}[1]{\textcolor[rgb]{0.38,0.63,0.69}{\textbf{\textit{#1}}}}
\newcommand{\AttributeTok}[1]{\textcolor[rgb]{0.49,0.56,0.16}{#1}}
\newcommand{\BaseNTok}[1]{\textcolor[rgb]{0.25,0.63,0.44}{#1}}
\newcommand{\BuiltInTok}[1]{\textcolor[rgb]{0.00,0.50,0.00}{#1}}
\newcommand{\CharTok}[1]{\textcolor[rgb]{0.25,0.44,0.63}{#1}}
\newcommand{\CommentTok}[1]{\textcolor[rgb]{0.38,0.63,0.69}{\textit{#1}}}
\newcommand{\CommentVarTok}[1]{\textcolor[rgb]{0.38,0.63,0.69}{\textbf{\textit{#1}}}}
\newcommand{\ConstantTok}[1]{\textcolor[rgb]{0.53,0.00,0.00}{#1}}
\newcommand{\ControlFlowTok}[1]{\textcolor[rgb]{0.00,0.44,0.13}{\textbf{#1}}}
\newcommand{\DataTypeTok}[1]{\textcolor[rgb]{0.56,0.13,0.00}{#1}}
\newcommand{\DecValTok}[1]{\textcolor[rgb]{0.25,0.63,0.44}{#1}}
\newcommand{\DocumentationTok}[1]{\textcolor[rgb]{0.73,0.13,0.13}{\textit{#1}}}
\newcommand{\ErrorTok}[1]{\textcolor[rgb]{1.00,0.00,0.00}{\textbf{#1}}}
\newcommand{\ExtensionTok}[1]{#1}
\newcommand{\FloatTok}[1]{\textcolor[rgb]{0.25,0.63,0.44}{#1}}
\newcommand{\FunctionTok}[1]{\textcolor[rgb]{0.02,0.16,0.49}{#1}}
\newcommand{\ImportTok}[1]{\textcolor[rgb]{0.00,0.50,0.00}{\textbf{#1}}}
\newcommand{\InformationTok}[1]{\textcolor[rgb]{0.38,0.63,0.69}{\textbf{\textit{#1}}}}
\newcommand{\KeywordTok}[1]{\textcolor[rgb]{0.00,0.44,0.13}{\textbf{#1}}}
\newcommand{\NormalTok}[1]{#1}
\newcommand{\OperatorTok}[1]{\textcolor[rgb]{0.40,0.40,0.40}{#1}}
\newcommand{\OtherTok}[1]{\textcolor[rgb]{0.00,0.44,0.13}{#1}}
\newcommand{\PreprocessorTok}[1]{\textcolor[rgb]{0.74,0.48,0.00}{#1}}
\newcommand{\RegionMarkerTok}[1]{#1}
\newcommand{\SpecialCharTok}[1]{\textcolor[rgb]{0.25,0.44,0.63}{#1}}
\newcommand{\SpecialStringTok}[1]{\textcolor[rgb]{0.73,0.40,0.53}{#1}}
\newcommand{\StringTok}[1]{\textcolor[rgb]{0.25,0.44,0.63}{#1}}
\newcommand{\VariableTok}[1]{\textcolor[rgb]{0.10,0.09,0.49}{#1}}
\newcommand{\VerbatimStringTok}[1]{\textcolor[rgb]{0.25,0.44,0.63}{#1}}
\newcommand{\WarningTok}[1]{\textcolor[rgb]{0.38,0.63,0.69}{\textbf{\textit{#1}}}}
\usepackage{longtable,booktabs,array}
\usepackage{calc} % for calculating minipage widths
% Correct order of tables after \paragraph or \subparagraph
\usepackage{etoolbox}
\makeatletter
\patchcmd\longtable{\par}{\if@noskipsec\mbox{}\fi\par}{}{}
\makeatother
% Allow footnotes in longtable head/foot
\IfFileExists{footnotehyper.sty}{\usepackage{footnotehyper}}{\usepackage{footnote}}
\makesavenoteenv{longtable}
\setlength{\emergencystretch}{3em} % prevent overfull lines
\providecommand{\tightlist}{%
  \setlength{\itemsep}{0pt}\setlength{\parskip}{0pt}}
% Algorithm Prime Triad v1.0.0 — pandoc header for pdfLaTeX.
\usepackage{newunicodechar}
\newunicodechar{≤}{\ensuremath{\le}}
\newunicodechar{≥}{\ensuremath{\ge}}
\newunicodechar{≠}{\ensuremath{\neq}}
\newunicodechar{Φ}{\ensuremath{\Phi}}
\newunicodechar{Γ}{\ensuremath{\Gamma}}
\newunicodechar{Δ}{\ensuremath{\Delta}}
\newunicodechar{Ω}{\ensuremath{\Omega}}
\newunicodechar{χ}{\ensuremath{\chi}}
\newunicodechar{⌊}{\ensuremath{\lfloor}}
\newunicodechar{⌋}{\ensuremath{\rfloor}}
\usepackage{microtype}
% Let long identifiers in code spans break after underscores.
\let\origunderscore\_
\renewcommand{\_}{\origunderscore\allowbreak}
\pdftrailerid{prime-algos-1.0.0}
\AtBeginDocument{\hypersetup{
  pdftitle={Algorithm Prime Triad: Candidate Intake, Adapter Descent, and Continuation Control},
  pdfauthor={Jeremy H. Carroll},
  pdfsubject={Version 1.0.0; 2026-09-26; preprint; DOI pending},
  pdfkeywords={algorithm prime triad, cancontinue predicate, candidate intake, adapter descent, Lean 4}
}}
\usepackage{bookmark}
\IfFileExists{xurl.sty}{\usepackage{xurl}}{} % add URL line breaks if available
\urlstyle{same}
\hypersetup{
  pdftitle={Algorithm Prime Triad: Candidate Intake, Adapter Descent, and Continuation Control},
  pdfauthor={Jeremy H. Carroll},
  colorlinks=true,
  linkcolor={Maroon},
  filecolor={Maroon},
  citecolor={Blue},
  urlcolor={Blue},
  pdfcreator={LaTeX via pandoc}}

\title{Algorithm Prime Triad: Candidate Intake, Adapter Descent, and
Continuation Control}
\author{Jeremy H. Carroll}
\date{Version 1.0.0 --- 2026-09-26}

\begin{document}
\makeatletter
\begin{center}
{\LARGE \@title\par}
\vspace{0.8em}
{\normalsize \@author\par}
\vspace{0.5em}
{\small \@date\par}
\end{center}
\makeatother

{
\hypersetup{linkcolor=}
\setcounter{tocdepth}{3}
\tableofcontents
}
\textbf{Series:} \texttt{008\_engine\_runtime} /
\texttt{000\_oe\_prime\_algorithms}\\
\textbf{Series position:} OE-000 (Part I --- Algorithm Prime Triad
Master Monograph)\\
\textbf{Author:} Jeremy H. Carroll\\
\textbf{Date:} 2026-09-26\\
\textbf{Version:} 1.0.0 (Unified Monograph combining Prime -1, 000, 001,
002, and CODA \(C^n\) Ascent)\\
\textbf{Document class:} Chapter / stem paper

\begin{center}\rule{0.5\linewidth}{0.5pt}\end{center}

\subsection{Why a Researcher Should
Care}\label{why-a-researcher-should-care}

Every recursive evaluator faces a fatal failure mode: it can always add
a dimension, rename an unresolved residual, or re-enter an evaluation
locus, thereby appearing to make progress forever without discharging
anything.

The \textbf{Algorithm Prime Triad (000, 001, 002)} establishes the
foundational operational control law of the Omega Engine:

\begin{enumerate}
\def\labelenumi{\arabic{enumi}.}
\tightlist
\item
  \textbf{Prime -1 \& 000:} Enforces strict locus grounding and a 7-step
  local packet declaration order
  (\(\text{Grounding} \to \text{Carrier} \to \text{Observer} \to \text{Primitive Difference} \to \text{ZeroRel} \to \text{Witness} \to \text{Quotient}\)).
\item
  \textbf{Prime 001:} Governs intake and inter-carrier transport
  soundness, enforcing that translation into an \texttt{OmegaPacket}
  does not confer validity and tracking all lost data in explicit loss
  ledgers (\(L_{\text{adapter}}\)).
\item
  \textbf{Prime 002:} Enforces the master continuation control law
  (\(\mathsf{CanContinue}_n\)) and six anti-cheat halting laws,
  guaranteeing finite, budgeted stopping.
\item
  \textbf{CODA \(C^n\) Ascent:} Formalizes recursive re-entry
  \(\Omega_{000}(L)\) at any locus without minting false terminal
  authority.
\end{enumerate}

Without a strict declaration, transport, and halting calculus, agentic
execution swarms degenerate into endless prompt retries, un-grounded
carrier hallucinations, and authority laundering.

\begin{center}\rule{0.5\linewidth}{0.5pt}\end{center}

\subsection{Refusal Box \& Scope Fence}\label{refusal-box-scope-fence}

\begin{Shaded}
\begin{Highlighting}[]
\NormalTok{No automatic conversion of raw prompts into valid cochains.}
\NormalTok{No degree, lift, or authority laundering.}
\NormalTok{No un{-}bounded recursive evaluation.}
\NormalTok{No un{-}grounded carrier hallucinations.}

\NormalTok{Candidate proposals under Prime {-}1 carry status: proposed\_only}
\NormalTok{and confer no validity until Prime 000 admission.}
\NormalTok{Step 10 terminalization emits LOCAL\_GREEN\_FOR\_SCOPE,}
\NormalTok{never an un{-}scoped global green certificate.}
\end{Highlighting}
\end{Shaded}

\begin{center}\rule{0.5\linewidth}{0.5pt}\end{center}

\subsection{Mathematical Foundations \&
Notation}\label{mathematical-foundations-notation}

The operational calculus operates over discrete finite simplicial
complexes \(K = (V, E, F)\) equipped with exact rational arithmetic over
\(\mathbb{Q}\):

\begin{enumerate}
\def\labelenumi{\arabic{enumi}.}
\tightlist
\item
  \textbf{Chains and Cochains:} \(C_0, C_1, C_2\) denote the rational
  vector spaces spanned by vertices, oriented edges, and oriented faces.
  Cochain spaces \(C^k = \operatorname{Hom}(C_k, \mathbb{Q})\) represent
  evaluation functions.
\item
  \textbf{Coboundary Operators:} The discrete derivative
  \(d_0: C^0 \to C^1\) assigns to potential \(\phi\) the edge difference
  \((d_0 \phi)(u, v) = \phi(v) - \phi(u)\). The curl operator
  \(d_1: C^1 \to C^2\) evaluates circulation around faces
  \((d_1 h)(u, v, w) = h(u, v) + h(v, w) - h(u, w)\), satisfying
  \(d_1 \circ d_0 = 0\).
\item
  \textbf{Exact Quotient Potential (\(\Phi_1\)):} Given a 1-cochain
  \(h \in C^1\), the quotient potential is the exact minimum L1 distance
  to the subspace of coboundaries (gauge transforms):
  \[\Phi_1(h) := \min_{\phi \in C^0} \|h - d_0 \phi\|_1\] This is
  evaluated exactly over \(\mathbb{Q}\) via rational linear programming,
  yielding a dual witness vector \(z \in [-1, 1]^{|E|}\) with
  \(d_0^* z = 0\) certifying the zero duality gap:
  \(\Phi_1(h) = \langle z, h \rangle\).
\item
  \textbf{Closure Charge (\(\Gamma_2\)):} The 2-obstruction measuring
  non-closure is: \[\Gamma_2(h) := \|d_1 h\|_1\]
\item
  \textbf{Discrete Hodge Decomposition:} For an inner product with
  positive edge weights \(W\), any 1-cochain decomposes orthogonally
  over \(\mathbb{Q}\) into exact, harmonic, and coexact components:
  \[h = d_0 \phi + h_{\text{harm}} + d_1^* \psi\] where
  \(h_{\text{harm}} \in \ker(d_1) \cap \ker(d_0^* W)\).
\end{enumerate}

\begin{center}\rule{0.5\linewidth}{0.5pt}\end{center}

\subsection{1. Prime -1 \& Prime 000: Candidate Proposal Engine \& Local
Declaration
Order}\label{prime--1-prime-000-candidate-proposal-engine-local-declaration-order}

\[\begin{aligned}
\text{Prime -1: Candidate Proposal} &\longrightarrow \text{Prime 001 Intake Check} \\
&\longrightarrow \text{Prime 000 Validation} \longrightarrow \text{DDA Execution}
\end{aligned}\]

Requiring an end-user to manually invent carriers, primitive
differences, or observers forces them to perform internal mathematical
tasks. \textbf{Prime -1 (Carrier Proposal Engine / Packet Compiler)}
acts as the front-end intake compiler:

\[\boxed{\begin{gathered}
\text{The Omega Engine does not begin by knowing the carrier.} \\
\text{It begins by proposing candidate carriers under scoped admission.}
\end{gathered}}\]

\subsubsection{\texorpdfstring{1.1 CandidatePacket \& IntakeTrace Schema
(\texttt{OE001-01})}{1.1 CandidatePacket \& IntakeTrace Schema (OE001-01)}}\label{candidatepacket-intaketrace-schema-oe001-01}

Prime -1 ingests a raw artifact and emits candidate
\texttt{CandidatePacket} structures. Pre-admission refusals live
strictly in an \texttt{IntakeTrace}
(\texttt{enters\_dda\_trace:\ false}) to prevent contaminating the
operational \texttt{DdaTrace}: \texttt{{[}DEFINITION{]}}
(\texttt{OE001-01})

\begin{Shaded}
\begin{Highlighting}[]
\FunctionTok{IntakeTrace}\KeywordTok{:}
\AttributeTok{  }\FunctionTok{raw\_artifact\_hash}\KeywordTok{:}\AttributeTok{ }\StringTok{"doc\_001\_sha256"}
\AttributeTok{  }\FunctionTok{candidate\_locus}\KeywordTok{:}\AttributeTok{ }\StringTok{"lines 1{-}12"}
\AttributeTok{  }\FunctionTok{candidate\_packet}\KeywordTok{:}
\AttributeTok{    }\FunctionTok{artifact\_id}\KeywordTok{:}\AttributeTok{ }\StringTok{"doc\_001\_ref"}
\AttributeTok{    }\FunctionTok{locus\_id}\KeywordTok{:}\AttributeTok{ }\StringTok{"locus\_frontmatter\_01"}
\AttributeTok{    }\FunctionTok{problem\_family}\KeywordTok{:}\AttributeTok{ }\StringTok{"claim | proof | narrative | code"}
\AttributeTok{    }\FunctionTok{primitive\_difference\_proposal}\KeywordTok{:}\AttributeTok{ }\StringTok{"metadata mismatch"}
\AttributeTok{    }\FunctionTok{carrier\_proposal}\KeywordTok{:}\AttributeTok{ }\StringTok{"document metadata block"}
\AttributeTok{    }\FunctionTok{zero\_rel\_proposal}\KeywordTok{:}\AttributeTok{ }\StringTok{"canonical metadata equality"}
\AttributeTok{    }\FunctionTok{observer\_proposal}\KeywordTok{:}\AttributeTok{ }\StringTok{"frontmatter schema checker"}
\AttributeTok{    }\FunctionTok{witness\_available}\KeywordTok{:}\AttributeTok{ }\CharTok{true}
\AttributeTok{    }\FunctionTok{quotient\_or\_refusal}\KeywordTok{:}\AttributeTok{ }\StringTok{"E or M"}
\AttributeTok{    }\FunctionTok{adapter\_confidence}\KeywordTok{:}\AttributeTok{ }\StringTok{"high"}
\AttributeTok{    }\FunctionTok{status}\KeywordTok{:}\AttributeTok{ }\StringTok{"proposed\_only"}
\AttributeTok{  }\FunctionTok{intake\_adapter\_status}\KeywordTok{:}\AttributeTok{ }\StringTok{"ADMITTED | REFUSED\_UNTYPED"}
\AttributeTok{  }\FunctionTok{enters\_dda\_trace}\KeywordTok{:}\AttributeTok{ }\CharTok{false}
\end{Highlighting}
\end{Shaded}

Prime 000 receives each candidate and emits one of seven formal
responses: \texttt{ACCEPT\_DECLARATION} ·
\texttt{REFUSE\_MISSING\_CARRIER} · \texttt{REFUSE\_MISSING\_OBSERVER} ·
\texttt{REFUSE\_MISSING\_WITNESS} · \texttt{REFUSE\_MISSING\_QUOTIENT} ·
\texttt{SPLIT\_PACKET} · \texttt{REQUEST\_LIFT}.

\subsubsection{\texorpdfstring{1.2 Prime 000 Local Packet Declaration
Order
(\texttt{OE001-02})}{1.2 Prime 000 Local Packet Declaration Order (OE001-02)}}\label{prime-000-local-packet-declaration-order-oe001-02}

Every admitted packet MUST be declared according to the 7-step local
runtime declaration order: \texttt{{[}DEFINITION{]}} (\texttt{OE001-02})

\begin{Shaded}
\begin{Highlighting}[]
\NormalTok{0. Grounding (cryptographic raw text / AST / commit hash anchor)}
\NormalTok{1. Carrier (simplicial complex / AST graph / discrete carrier)}
\NormalTok{2. Observer packet / detector (resolution, probe, gauge)}
\NormalTok{3. Primitive difference (cochain h in C1 or defect functional)}
\NormalTok{4. ZeroRel rule (gauge equivalence / coboundary relation)}
\NormalTok{5. Witness or refusal route (replay certificate chi or typed refusal)}
\NormalTok{6. Optional quotient, degree, lift, or terminal policy}
\end{Highlighting}
\end{Shaded}

\subsubsection{\texorpdfstring{1.3 Layered Multi-Carrier Lenses \&
LensSpan Model
(\texttt{OE001-03})}{1.3 Layered Multi-Carrier Lenses \& LensSpan Model (OE001-03)}}\label{layered-multi-carrier-lenses-lensspan-model-oe001-03}

A complex artifact (paper, codebase, financial plan) is \textbf{not} a
single monolithic cochain complex. It is a \textbf{bundle of carriers}
extracted via layered partial spans: \texttt{{[}SCOPE\ FENCE{]}}
(\texttt{OE001-03})

\[A \xleftarrow{s_L} S_L \xrightarrow{p_L} K_L\]

where \(S_L\) is the selected source span, \(K_L\) is the extracted
carrier, and span composition proceeds by pullback in the category
\(\mathbf{LensSpan}(A)\) under canonical inclusion subcomplexes:

\[\begin{aligned}
L_0 \ (\text{lexical}) &\longrightarrow L_1 \ (\text{syntactic}) \longrightarrow L_2 \ (\text{semantic}) \\
&\longrightarrow L_3 \ (\text{proof}) \longrightarrow L_4 \ (\text{rhetoric}) \longrightarrow L_5 \ (\text{strategic})
\end{aligned}\]

\subsubsection{\texorpdfstring{1.4 Scoped Terminalization
(\texttt{OE001-04})}{1.4 Scoped Terminalization (OE001-04)}}\label{scoped-terminalization-oe001-04}

Terminalization (\(T\)) emits scoped local status rather than un-scoped
global green: \texttt{{[}OPERATIONAL\ RULE{]}} (\texttt{OE001-04})

\[\begin{aligned}
\text{Step 10: } T &\longrightarrow \text{Terminalize local repair pass:} \\
&\quad \mathsf{LOCAL\_GREEN\_FOR\_SCOPE}(\text{scope}, \text{remaining\_obligations})
\end{aligned}\]

Terminalization evaluates the active cochain state and routes through
two distinct terminal paths:

\begin{enumerate}
\def\labelenumi{\arabic{enumi}.}
\tightlist
\item
  \textbf{Zero-Residual Terminalization (\(T_{\text{res}}\)):} If the
  quotient residual vanishes (\(\Phi_1 = 0\)) and the closure charge is
  exact (\(\Gamma_2 = 0\)) under an adequate detector, the pass seals
  with \(\mathsf{LOCAL\_GREEN\_FOR\_SCOPE}\). The emitted receipt
  registers the exact potential and witnesses, certifying completion
  strictly within the declared carrier scope.
\item
  \textbf{Terminal Failure / Refusal (\(F_{\text{term}}\)):} If an open
  residual (\(\Phi_1 > 0\)) or closure failure (\(\Gamma_2 > 0\))
  persists, terminalization is refused. The evaluator halts with typed
  terminal failure \(F_{\text{term}}\) (\texttt{REFUSED\_OPEN\_RESIDUAL}
  or \texttt{REFUSED\_BLIND\_OBSERVER}), preventing un-repaired defects
  from being sealed as green.
\end{enumerate}

Under no circumstances does terminalization convert a local repair into
an un-scoped global green certificate.

\begin{center}\rule{0.5\linewidth}{0.5pt}\end{center}

\subsection{2. Prime 001: Adapter Descent Soundness \& Transport
Verification}\label{prime-001-adapter-descent-soundness-transport-verification}

\subsubsection{\texorpdfstring{2.1 Dual Prime 001 Placement Protocol
(\texttt{OE002-01})}{2.1 Dual Prime 001 Placement Protocol (OE002-01)}}\label{dual-prime-001-placement-protocol-oe002-01}

To prevent un-grounded or invalid translations from contaminating the
execution log, Prime 001 is deployed at \textbf{two distinct boundary
points}: \texttt{{[}OPERATIONAL\ RULE{]}} (\texttt{OE002-01})

\[\begin{aligned}
\text{Prime (-1)} &\longrightarrow \text{Prime 001 [Intake Adapter]} \longrightarrow \text{Prime 000 [Admit]} \\
&\longrightarrow \text{DDA Runner} \longrightarrow \text{Prime 001 [Transport Adapter]} \longrightarrow \text{Prime 002 [Control]}
\end{aligned}\]

\begin{enumerate}
\def\labelenumi{\arabic{enumi}.}
\tightlist
\item
  \textbf{Intake Adapter Check:} Evaluates raw artifact \(\to\)
  candidate packet translation before Prime 000 admission.
\item
  \textbf{Transport Adapter Check:} Evaluates inter-carrier transport
  steps whenever a DDA operation changes carrier.
\end{enumerate}

\subsubsection{\texorpdfstring{2.2 Adapter Descent Soundness Rule
(\texttt{OE002-02})}{2.2 Adapter Descent Soundness Rule (OE002-02)}}\label{adapter-descent-soundness-rule-oe002-02}

Translation of a native claim into an \texttt{OmegaPacket} is a
representation mapping, not a proof of validity. An intake or transport
adapter \(\alpha: C_{\text{source}} \to C_{\text{target}}\) is sound if
and only if it preserves operations, equations, witnesses, and authority
fences while recording all lost data in an explicit loss ledger
\(L_{\text{adapter}}\): \texttt{{[}DEFINITION{]}} (\texttt{OE002-02})

\begin{Shaded}
\begin{Highlighting}[]
\FunctionTok{AdapterContract}\KeywordTok{:}
\AttributeTok{  }\FunctionTok{source\_carrier}\KeywordTok{:}\AttributeTok{ C\_source}
\AttributeTok{  }\FunctionTok{target\_carrier}\KeywordTok{:}\AttributeTok{ C\_target}
\AttributeTok{  }\FunctionTok{preserved\_operations}\KeywordTok{:}\AttributeTok{ }\KeywordTok{[}\AttributeTok{d0}\KeywordTok{,}\AttributeTok{ d1}\KeywordTok{]}
\AttributeTok{  }\FunctionTok{preserved\_equations}\KeywordTok{:}\AttributeTok{ }\KeywordTok{[}\AttributeTok{d1\_d0\_zero}\KeywordTok{]}
\AttributeTok{  }\FunctionTok{preserved\_witnesses}\KeywordTok{:}\AttributeTok{ }\KeywordTok{[}\AttributeTok{dual\_certificate}\KeywordTok{]}
\AttributeTok{  }\FunctionTok{loss\_ledger}\KeywordTok{:}\AttributeTok{ L\_adapter}
\end{Highlighting}
\end{Shaded}

If \(L_{\text{adapter}} \neq \varnothing\), the adapter emits
\texttt{ADAPTER\_DEBT} and halts automatic promotion.

\subsubsection{\texorpdfstring{2.3 Refusal \& Debt Taxonomy
(\texttt{OE002-03})}{2.3 Refusal \& Debt Taxonomy (OE002-03)}}\label{refusal-debt-taxonomy-oe002-03}

\begin{enumerate}
\def\labelenumi{\arabic{enumi}.}
\tightlist
\item
  \texttt{REFUSED\_UNTYPED}: Raw input cannot be mapped to any known
  carrier signature or observer probe.
\item
  \texttt{ADAPTER\_INADEQUATE}: Candidate carrier mapping exists but
  fails boundary operator or witness preservation.
\item
  \texttt{REFUSED\_ADAPTER\_DEBT}: Transport across carrier boundaries
  failed adapter soundness checks (\texttt{C\ ·\ M\ ·\ T}).
\end{enumerate}

\begin{center}\rule{0.5\linewidth}{0.5pt}\end{center}

\subsection{3. Prime 002: Master Continuation Predicate \& Anti-Cheat
Control}\label{prime-002-master-continuation-predicate-anti-cheat-control}

\subsubsection{\texorpdfstring{3.1 Master Control Loop \& Predicate
\(\mathsf{CanContinue}_n\)
(\texttt{OE003-01})}{3.1 Master Control Loop \& Predicate \textbackslash mathsf\{CanContinue\}\_n (OE003-01)}}\label{master-control-loop-predicate-mathsfcancontinue_n-oe003-01}

At stage \(n\), continuation to stage \(n+1\) is lawful if and only if
predicate \(\mathsf{CanContinue}_n\) evaluates to \texttt{True}:
\texttt{{[}DEFINITION{]}} (\texttt{OE003-01})

\[\begin{aligned}
\mathsf{CanContinue}_n :\iff &\mathsf{OutstandingObligation}(r_n) \wedge \mathsf{HigherObligationWitness}(r_n) \wedge \\
&\mathsf{NextCarrierDeclared}(r_n) \wedge \mathsf{AdapterSound}_{n\to n+1} \wedge \\
&\mathsf{DetectorAdequate}_{n+1} \wedge \mathsf{MinimalLicensedLift}_{n\to n+1} \wedge \\
&\mathsf{SkeletonPreserved}_{n\to n+1} \wedge \mathsf{BudgetAvailable}_n \wedge \\
&\neg\mathsf{GTLAContinuationRefusal}_n
\end{aligned}\]

The three-condition form groups these conjuncts into structural
validity, resource availability, and absence of terminal refusal. The
Lean theorem \texttt{Prime.cc9\_implies\_cc3} formalizes this equivalent
regrouping.

\subsubsection{\texorpdfstring{3.2 Six Anti-Cheat Halting Laws
(\texttt{OE003-02})}{3.2 Six Anti-Cheat Halting Laws (OE003-02)}}\label{six-anti-cheat-halting-laws-oe003-02}

\begin{quote}
\textbf{Halting Law Constraints:}

\begin{enumerate}
\def\labelenumi{\arabic{enumi}.}
\tightlist
\item
  \textbf{No Degree Laundering:} Adding \(C^{n+1}\) requires an
  exhibited \((n+1)\)-dimensional obligation.
\item
  \textbf{No Lift Laundering:} A representation extension must prove no
  same-carrier repair exists.
\item
  \textbf{No Anomalon Evasion:} A stable residual cannot be perpetually
  renamed to avoid terminalization.
\item
  \textbf{Old-Skeleton Preservation:} A continuation must preserve
  certified data or emit a loss ledger.
\item
  \textbf{Detector Revalidation:} Repacketization must revalidate
  detector adequacy at degree \(n+1\).
\item
  \textbf{First-Unresolved-Degree Honesty:} If \(n_0\) is unresolved,
  claims above \(n_0\) are conditional: \(\mathsf{CoherentUpTo}(n_0-1)\)
  and \(\mathsf{ResidualAt}(n_0)\).
\end{enumerate}
\end{quote}

\subsubsection{\texorpdfstring{3.3 Bounded Operational Termination
(\texttt{OE003-03})}{3.3 Bounded Operational Termination (OE003-03)}}\label{bounded-operational-termination-oe003-03}

Given a finite non-replenished budget \(B_{\text{tot}} \in \mathbb{N}\)
(or scaled rational budget) where every counted operation debits the
shared budget by at least a strictly positive cost floor
\(\delta_{\text{control}} > 0\) and no operation can begin without
sufficient remaining budget (\(B \ge \delta_{\text{control}}\)), the
number of executed operations is bounded by:
\texttt{{[}CONDITIONAL\ THEOREM{]}} (\texttt{OE003-03})

\[\#\mathrm{ops} \le \left\lfloor \frac{B_{\text{tot}}}{\delta_{\text{control}}} \right\rfloor\]

This bound (\texttt{Prime.ops\_le\_floor\_div}) formalizes the
natural-number arithmetic consequence of the non-replenished debit
invariant (\(n \cdot \delta \le B\)). Total runtime termination
additionally assumes that each individual operational checker halts in
finite time. If budgets are tracked in exact rationals \(\mathbb{Q}\),
the bound applies via scaling by the least common denominator.

\subsubsection{\texorpdfstring{3.4 Local Queue Scheduler
(\texttt{OE003-04}) \& CODA Taxonomy
(\texttt{OE003-05})}{3.4 Local Queue Scheduler (OE003-04) \& CODA Taxonomy (OE003-05)}}\label{local-queue-scheduler-oe003-04-coda-taxonomy-oe003-05}

For local execution on a fixed carrier, Prime 002 acts as a finite
residual queue scheduler. To prevent anomalon evasion (Law 3), progress
requires strict lexicographical descent on the diagnostic tuple:
\texttt{{[}OPERATIONAL\ RULE{]}} (\texttt{OE003-04})

\[\text{Descent} :\iff (\Gamma_{2, n+1}, \Phi_{1, n+1}, \text{defect\_count}_{n+1}) <_{\text{lex}} (\Gamma_{2, n}, \Phi_{1, n}, \text{defect\_count}_n)\]

Witnessed exploration or residual fingerprint changes represent
exploratory transitions, not mathematical progress. Across steps on a
fixed carrier, continuation requires quantitative potential descent
\(\Delta \Phi_1 \le 0\) evaluated using the exact quotient norm
\(\|[h]\|_{Q1}\). Under a licensed lift across carriers
(\(C^n \to C^{n+1}\)), \(\Phi_1\) may temporarily increase if the
representation extension exposes higher-dimensional boundary
obligations, provided the lift license explicitly certifies that no
same-carrier descent is possible: \texttt{{[}FINITE\ THEOREM{]}}
(\texttt{OE003-06}).

\paragraph{\texorpdfstring{CODA Pipeline-Order Taxonomy
(\texttt{OE003-05})}{CODA Pipeline-Order Taxonomy (OE003-05)}}\label{coda-pipeline-order-taxonomy-oe003-05}

CODA formalizes local obstruction analysis as a strict 7-stage pipeline
sequence, with \(Q\) as The Continuum: \texttt{{[}DEFINITION{]}}
(\texttt{OE003-05})

\[\begin{aligned}
G \ (\text{Grounding}) &\longrightarrow C \ (\text{Carrier}) \longrightarrow J \ (\text{Observer}) \\
&\longrightarrow \Delta \ (\text{Primitive Difference}) \longrightarrow Z \ (\text{ZeroRel}) \\
&\longrightarrow \chi \ (\text{Witness}) \longrightarrow Q \ (\text{Quotient})
\end{aligned}\]

\begin{enumerate}
\def\labelenumi{\arabic{enumi}.}
\tightlist
\item
  \textbf{Grounding (\(G\)):} Raw artifact cryptographic digest and
  origin anchor.
\item
  \textbf{Carrier (\(C\)):} Explicit simplicial complex or discrete
  chain module representing the substrate.
\item
  \textbf{Observer (\(J\)):} Finite projection and resolution window
  specifying measurement capacity.
\item
  \textbf{Primitive Difference (\(\Delta\)):} Concrete difference
  cochain \(h \in C^1\) between competing models or steps.
\item
  \textbf{ZeroRel (\(Z\)):} Equivalence relation specifying exact
  triviality or gauge equivalence (\(\operatorname{im} d_0\)).
\item
  \textbf{Witness (\(\chi\)):} Explicit dual witness certificate
  bounding the duality gap or confirming exact repair.
\item
  \textbf{Quotient (\(Q\)):} The Continuum, the quotient space
  \(C^1 / \operatorname{im}(d_0)\) hosting the residual obstruction.
\end{enumerate}

Output status taxonomy:

\begin{itemize}
\tightlist
\item
  \texttt{LOCAL\_GREEN\_FOR\_SCOPE}: Exact zero residual
  (\(\Phi_1 = 0, \Gamma_2 = 0\)) under an adequate observer.
\item
  \texttt{REFUSED\_UNTYPED}: Candidate locus lacking a declared carrier
  or grounding anchor.
\item
  \texttt{ADAPTER\_INADEQUATE}: Loss ledger identifies loss of essential
  invariants or unpreservable equations.
\item
  \texttt{HALTED\_BUDGET\_EXHAUSTED}: Operational budget depleted
  (\(B < \delta_{\text{control}}\)).
\item
  \texttt{HALTED\_STAGNANT\_PROGRESS}: Lookback window detects lack of
  strictly decreasing residual mass.
\item
  \texttt{GTLA\_STOP}: Terminal local stop halted with an unresolved
  obstruction that cannot be discharged under available resources.
\item
  \texttt{GTTA\_GLOBAL\_RESIDUAL}: Ambient structural obstruction
  surviving into all admissible representations.
\end{itemize}

\begin{center}\rule{0.5\linewidth}{0.5pt}\end{center}

\subsection{\texorpdfstring{4. CODA Recursive Re-Entry and \(C^n\)
Ascent
Protocol}{4. CODA Recursive Re-Entry and C\^{}n Ascent Protocol}}\label{coda-recursive-re-entry-and-cn-ascent-protocol}

\subsubsection{\texorpdfstring{4.1 CODA Recursive Re-Entry Protocol
(\texttt{OE004-01})}{4.1 CODA Recursive Re-Entry Protocol (OE004-01)}}\label{coda-recursive-re-entry-protocol-oe004-01}

Algorithm Prime 000 may be recursively re-entered at any CODA locus
\(L \in \text{CODA}\) (object, morphism, witness, residual, quotient,
adapter, or terminal claim).

A \textbf{locus} \(L\) is the entity under examination, whereas a
\textbf{carrier} \(C_L\) is the declared local mathematical structure in
which comparisons and primitive differences on \(L\) are represented. In
general, \(L \neq C_L\) (identifying them without an explicit typing
witness silently promotes an arbitrary locus into a carrier, bypassing
the Prime -1 proposal engine and Prime 000 declaration order).

The recursive re-entry pipeline requires proposing and declaring a
candidate carrier \(C_L\) for locus \(L\) before running Prime 000 on
\((L, C_L)\): \texttt{{[}PROTOCOL\ SPECIFICATION{]}} (\texttt{OE004-01})

\[\boxed{\text{locus } L \longrightarrow \text{propose/declare carrier } C_L \longrightarrow \text{run Prime 000 on } (L, C_L)}\]

\[\begin{aligned}
\Omega_{000}(L) \Longrightarrow \bigl(&\text{Carrier}=C_L,\ \mathrm{ZeroRel}_L,\ \text{Observer}=J_L, \\
&\text{Witness}=\chi_L,\ \text{Residual}=R_L,\ \text{Quotient}=Q_L \bigr)
\end{aligned}\]

subject to standard Prime 000 admission requirements.

\subsubsection{\texorpdfstring{4.2 Lawful \(C^n\) Ascent Protocol
(\texttt{OE004-02})}{4.2 Lawful C\^{}n Ascent Protocol (OE004-02)}}\label{lawful-cn-ascent-protocol-oe004-02}

Degree ascent \(C^k \longrightarrow C^{k+1}\) is permitted if and only
if the paper declares: \texttt{{[}DEFINITION{]}} (\texttt{OE004-02})

\begin{enumerate}
\def\labelenumi{\arabic{enumi}.}
\tightlist
\item
  Carrier complex \(C^k\), Differential \(d^k: C^k \to C^{k+1}\), Next
  Differential \(d^{k+1}: C^{k+1} \to C^{k+2}\), and exact nilpotency
  \(d^{k+1} \circ d^k = 0\).
\item
  \textbf{Monotone Phase Embedding Definition:} Given an assumed
  monotone map embedding admission into execution, if admission precedes
  measurement
  (\(\operatorname{admit}(P) \le \operatorname{measure}(P)\)), then the
  order-preserving structure guarantees that admission precedence
  reflects in execution for any packet pair where the map is monotone
  (\texttt{Prime.embedding\_order\_preserving}). This formalization
  proves the property of the assumed monotone structure; constructing a
  concrete causal trace embedding and differential nilpotency
  (\(d \circ d = 0\)) are separate, independent verification
  obligations.
\end{enumerate}

\subsubsection{\texorpdfstring{4.3 Stop Conditions: GTLA Stop vs.~CODA-A
Global Residual
(\texttt{OE004-03})}{4.3 Stop Conditions: GTLA Stop vs.~CODA-A Global Residual (OE004-03)}}\label{stop-conditions-gtla-stop-vs.-coda-a-global-residual-oe004-03}

\begin{itemize}
\tightlist
\item
  \textbf{\texttt{GTLA\_STOP}:} Local authority construction halted due
  to missing carrier, detector, witness, quotient, or exhausted
  resources.
\item
  \textbf{\texttt{GTTA\_GLOBAL\_RESIDUAL}:} Ambient structural
  obstruction \(R_U\) surviving transport into all admissible adapters.
\end{itemize}

\subsubsection{\texorpdfstring{4.4 Observer Envelopes \& Evaluator
Defect Linearity (\texttt{OE004-04}, \texttt{OE004-05},
\texttt{OE004-06})}{4.4 Observer Envelopes \& Evaluator Defect Linearity (OE004-04, OE004-05, OE004-06)}}\label{observer-envelopes-evaluator-defect-linearity-oe004-04-oe004-05-oe004-06}

\paragraph{\texorpdfstring{Observer Envelope Comparability
(\texttt{OE004-04})}{Observer Envelope Comparability (OE004-04)}}\label{observer-envelope-comparability-oe004-04}

Observer envelopes \(J_1, J_2\) are comparable only under shared
grounding and explicit coordinate transformation bounds:
\texttt{{[}SCOPE\ FENCE{]}} (\texttt{OE004-04}). Comparing evaluation
scores across ungrounded or incomparable observers is strictly refused.

\paragraph{\texorpdfstring{Evaluator Defect Linearity
(\texttt{OE004-05})}{Evaluator Defect Linearity (OE004-05)}}\label{evaluator-defect-linearity-oe004-05}

Let \(\Omega\) be a finite probability space with joint probability
distribution \(\mathbb{P}\). Let \(E_g: A \otimes B \to [0, 1]\) be a
global composite evaluator and \(E_1: A \to [0, 1], E_2: B \to [0, 1]\)
be local component evaluators. Define the global pass rate and joint
local pass rate: \texttt{{[}FINITE\ THEOREM{]}} (\texttt{OE004-05})

\[\mathrm{GPR} := \mathbb{E}[E_g(A \otimes B)], \qquad \mathrm{LPR} := \mathbb{E}[E_1(A) E_2(B)]\]

The pointwise evaluator defect is defined as:
\[\Delta_E(A, B) := E_g(A \otimes B) - E_1(A) E_2(B)\] By linearity of
expectation over the joint probability law:
\[\mathrm{GPR} - \mathrm{LPR} = \mathbb{E}[E_g(A \otimes B) - E_1(A) E_2(B)] = \mathbb{E}[\Delta_E]\]
When evaluators decouple without cross-module obstruction
(\(\Delta_E = 0\) pointwise), local pass rates predict global pass rates
identically (\(\mathrm{GPR} = \mathrm{LPR}\)). Note that zero expected
defect (\(\mathbb{E}[\Delta_E] = 0\)) can also arise via cancellation
unless \(\Delta_E \ge 0\); only a non-zero expected defect forces a
non-zero defect on some supported instances.

\paragraph{\texorpdfstring{Six-Role EPK Overlay Diagnostic Scope Fence
(\texttt{OE004-06})}{Six-Role EPK Overlay Diagnostic Scope Fence (OE004-06)}}\label{six-role-epk-overlay-diagnostic-scope-fence-oe004-06}

The six-role EPK overlay enforces strict role boundaries:
\texttt{{[}SCOPE\ FENCE{]}} (\texttt{OE004-06}). Evaluators, improvers,
adversaries, certifiers, archivers, and governors evaluate strictly
within their declared diagnostic scopes.

Engine self-examination (\(\Omega_{000}(L_{\text{self}})\)) and its
anti-cheat queue upgrade rules are specified in
\texttt{docs/residual\_aware\_can\_continue.md}:
\texttt{{[}MODEL\ CONSEQUENCE{]}} (\texttt{OE004-07},
\texttt{OE004-08}).

\begin{center}\rule{0.5\linewidth}{0.5pt}\end{center}

\subsection{5. Summary \& Claim-ID Index}\label{summary-claim-id-index}

\begin{longtable}[]{@{}
  >{\raggedright\arraybackslash}p{(\linewidth - 6\tabcolsep) * \real{0.2500}}
  >{\raggedright\arraybackslash}p{(\linewidth - 6\tabcolsep) * \real{0.2500}}
  >{\raggedright\arraybackslash}p{(\linewidth - 6\tabcolsep) * \real{0.2500}}
  >{\raggedright\arraybackslash}p{(\linewidth - 6\tabcolsep) * \real{0.2500}}@{}}
\toprule\noalign{}
\begin{minipage}[b]{\linewidth}\raggedright
Claim ID
\end{minipage} & \begin{minipage}[b]{\linewidth}\raggedright
Type
\end{minipage} & \begin{minipage}[b]{\linewidth}\raggedright
Section
\end{minipage} & \begin{minipage}[b]{\linewidth}\raggedright
Summary
\end{minipage} \\
\midrule\noalign{}
\endhead
\bottomrule\noalign{}
\endlastfoot
\texttt{OE001-01} & \texttt{Definition} & §1.1 & CandidatePacket \&
IntakeTrace schema definition \\
\texttt{OE001-02} & \texttt{Definition} & §1.2 & Prime 000 7-step local
packet declaration order \\
\texttt{OE001-03} & \texttt{Scope} & §1.3 & Multi-carrier lens span
pullback model \\
\texttt{OE001-04} & \texttt{Model\ consequence} & §1.4 & Scoped
terminalization \texttt{LOCAL\_GREEN\_FOR\_SCOPE} \\
\texttt{OE002-01} & \texttt{Scope} & §2.1 & Dual Prime 001 boundary
placement (Intake \& Transport) \\
\texttt{OE002-02} & \texttt{Definition} & §2.2 & Adapter contract
definition \& loss ledger structure \\
\texttt{OE002-03} & \texttt{Definition} & §2.3 & Refusal \& debt
taxonomy (\texttt{REFUSED\_UNTYPED}, \texttt{ADAPTER\_INADEQUATE}) \\
\texttt{OE003-01} & \texttt{Definition} & §3.1 & Continuation predicate
definition \& equivalent 3-form regrouping \\
\texttt{OE003-02} & \texttt{Protocol\ specification} & §3.2 & Six
anti-cheat halting laws \\
\texttt{OE003-03} & \texttt{Conditional\ theorem} & §3.3 & Bounded
operational termination law \\
\texttt{OE003-04} & \texttt{Operational\ rule} & §3.4 & Local queue
scheduler \& strict quotient descent condition \\
\texttt{OE003-05} & \texttt{Definition} & §3.4 & CODA taxonomy in
pipeline order with \(Q\) as The Continuum \\
\texttt{OE003-06} & \texttt{Finite\ theorem} &
\texttt{docs/residual\_aware\_can\_continue.md} & Residual-aware
\(\mathsf{CanContinue}_n\) with quantitative \(\Delta \Phi \le 0\)
bound \\
\texttt{OE004-01} & \texttt{Scope} & §4.1 & CODA recursive re-entry
protocol \(\Omega_{000}(L)\) over declared carrier \(C_L\)
(\(L \neq C_L\)) \\
\texttt{OE004-02} & \texttt{Definition} & §4.2 & Lawful ascent protocol
\& assumed monotone phase embedding \\
\texttt{OE004-03} & \texttt{Definition} & §4.3 & \texttt{GTLA\_STOP}
vs.~\texttt{GTTA\_GLOBAL\_RESIDUAL} stop conditions \\
\texttt{OE004-04} & \texttt{Scope} & §4.4 & Observer envelope
comparability \\
\texttt{OE004-05} & \texttt{Finite\ theorem} & §4.4 & Evaluator defect
linearity \(\mathrm{GPR} - \mathrm{LPR} = \mathbb{E}[\Delta_E]\) \\
\texttt{OE004-06} & \texttt{Scope} & §4.4 & Six-role EPK overlay
diagnostic scope fence \\
\texttt{OE004-07} & \texttt{Model\ consequence} &
\texttt{docs/residual\_aware\_can\_continue.md} & Self-examination
re-entry path \(\Omega_{000}(L_{\text{self}})\) for engine state \\
\texttt{OE004-08} & \texttt{Model\ consequence} &
\texttt{docs/residual\_aware\_can\_continue.md} & Anti-cheat laws for
self-proposal \& progress queue upgrade \\
\end{longtable}

\begin{center}\rule{0.5\linewidth}{0.5pt}\end{center}

\subsection{6. References}\label{references}

\begin{enumerate}
\def\labelenumi{\arabic{enumi}.}
\tightlist
\item
  J. H. Carroll. \emph{Discrete Deformation Calculus: A Finite Trace
  Calculus for Obstructions, Exact Repair, and Bounded Replay}, v1.0.0
  (2026).
\item
  J. H. Carroll. \emph{Local Green, Global Red I: Evaluated Path Kernels
  for Composed Agent Reliability}, v1.2.3 (2026). Zenodo, DOI:
  10.5281/zenodo.22881911.
\item
  J. H. Carroll. \emph{CODA-FOC: Finite Obstruction Calculus --- Exact
  Quotient Diagnostics, Dual Witnesses and Two-Stage Coherence}. Omega
  Engine Research Programme (2026).
\item
  J. H. Carroll. \emph{CODA-GTLA: General Theory of Layered Authority
  --- A Leftover Class is Not Standing}. Omega Engine Research Programme
  (2026).
\item
  The Lean 4 Theorem Prover, v4.34.0. \texttt{https://lean-lang.org}.
\end{enumerate}

\end{document}
